% ------------------------------------------------------------
% Page 95
% ------------------------------------------------------------

\textbf{Le Temple - Ecole des Oubrets}

Vers 1840, 50 ans après la Révolution et la fin des persécutions des protestants,
les églises réformées se sont peu à peu reconstituées. Le temple de Meyrueis a été
rebâti, après avoir été rasé lors de La Révocation de l’Edit de Nantes en 1685, non
plus au centre de la ville mais un peu à l’extérieur, sur l’emplacement d’une ancienne
glacière. La commune de Gatuzieres édifie également un petit temple et les habitants
des Oubrets, hameau situé à 11 km de Meyrueis, pour ne pas être en reste, sont
autorisés à aménager un lieu de culte.

\begin{quote}
\textit{«... Ce temple est fort original car il a une double fonction : en semaine il
abrite l’école communale du hameau, et deux dimanches par mois, ainsi qu’aux fêtes,
il retrouve sa fonction cultuelle.}

\textit{Cette situation se maintiendra jusque dans les années 1950, époque de la fermeture
de l’école et de la suppression des cultes aux Oubrets.}

\textit{Des témoignages oraux recueillis dans les années 1975 font état d’une pratique
plutôt curieuse, ayant cours au début de notre siècle. Le temple des Oubrets étant
dépourvu de cloche, les hommes du hameau, tous chasseurs de père en fils, se
rendaient, un moment avant le culte, sur un rocher dominant le village afin de tirer
des coups de fusil en l’air, pour rappeler aux habitants du lieu et des mas voisins de
Valbelle, Campredon et Campis, la tenue du service...}

\textit{Il faut noter aussi que l’unique famille catholique du hameau participait assidûment
aux cultes et aux veillées organisées dans ce local. »}
\end{quote}

\begin{flushright}
Philippe Chambon\\
\textit{« Itinéraires Protestants en Languedoc : Les Cévennes »} p.78
\end{flushright}

André Bazalgette, Instituteur aux Oubrets indique, dans ses souvenirs, que l’école fut
définitivement fermée en 1944 et que le temple, désaffecté quelque temps après, fut acquis
par les frères Rousset qui en firent leur maison d’habitation. Lorsque Adrien Rousset se
retrouva seul il s’installa à Pourcarès et plus tard, après son décès, l’ancien temple - école fut
acheté par Gilbert Martin boucher à Florac, fils de Gabriel Martin de Pourcarès.
